\begin{table}[ht]
\centering
{\tablefont
\resizebox{\textwidth}{!}{%
\begin{tabular}{rrrrll}
\toprule
Tier & Targets & $L_0$ runtime & GREG runtime & IPF runtime & Main outcome \\
\midrule
1 & 250 & \tbc & \tbc & \tbc & \tbc \\
2 & 500 & \tbc & \tbc & \tbc & \tbc \\
3 & 1{,}000 & \tbc & \tbc & \tbc & \tbc \\
4 & 2{,}500 & \tbc & \tbc & \tbc & \tbc \\
5 & 5{,}000 & \tbc & \tbc & \tbc & \tbc \\
6 & 10{,}000 & \tbc & \tbc & \tbc & \tbc \\
7 & \tbc[full selected tier] & \tbc & \tbc & \tbc & \tbc \\
\bottomrule
\end{tabular}
}
}
\caption{Tier 2 scaling frontier. Each row uses the same cloned unit universe and increases the
number of selected targets. Runtime cells should be replaced by failure labels when a method does
not complete.}
\label{tab:scaling_frontier}
\tablenote{The final paper may split this table into separate runtime and failure-mode tables if
the notes column becomes too dense.}
\end{table}
